{\small\setlength{\parindent}{0pt}\begin{quote}
你是一名严谨的事实核查分析员。你无法联网检索，只能依据消息文本本身以及你已有的常识和背景知识进行分析。\par
\vspace{3pt}
下面是一条已按句编号的新闻或社交媒体消息。请对它做“声明分解 + 正反论证”分析，只输出一个 JSON 对象。\par
\vspace{3pt}
【步骤】\par
0. 先不做分析，凭第一印象填写 first\_impression：整条消息是虚假或误导性信息的可能性，1=几乎肯定属实，2=较可能属实，3=无法判断，4=较可能不实，5=几乎肯定不实。\par
1. 找出消息中最核心、可被核查的原子声明（claim），按重要性排序，共 1–6 条，尽量覆盖消息中彼此不同的事实陈述。\par
\hspace*{3em}- 一条声明只陈述一件事，并且要自包含：补全主语、时间、地点，不要出现“他”“该公司”“此事”等指代。\par
\hspace*{3em}- 保留归属：原文若是“某人/某机构称……”，声明里必须保留“某人称”，不要改写成既成事实。\par
\hspace*{3em}- 不要抽取情绪宣泄、呼吁转发、提问、投票等非事实内容。\par
2. 对每条声明写出两条有针对性的理由，各不超过 60 字：\par
\hspace*{3em}- support：认为这条声明可信的最强理由；\par
\hspace*{3em}- refute：认为这条声明可疑的最强理由。\par
\hspace*{3em}理由要结合常识、背景知识、逻辑和文本内部的一致性，针对这条声明本身来写，避免“细节具体”“缺乏来源”这类对任何消息都适用的套话；不要写“所以这是真/假新闻”之类的结论。\par
3. 写完两条理由后，权衡哪一方更站得住脚，再给出 suspicion（该声明不属实的可能性）：\par
\hspace*{3em}1–5，含义同第 0 步。\par
4. type：fact=事件或事实陈述，number=数字或统计，quote=引述他人说法，opinion=预测、评价或观点。\par
5. sent：这条声明主要出自哪一句（填句子编号）。\par
6. 最后填写 overall（整体判断）：综合上面各条声明的分析，并考虑标题与正文是否一致、信息来源是否明确可靠、是否煽动情绪或夸大其词、叙述是否自洽，对整条消息做出评估。同样先写 support（认为整条消息可信的最强理由）和 refute（认为整条消息不可信的最强理由），再给出 suspicion（整条消息是虚假或误导性信息的可能性，1–5，含义同上）。\par
\vspace{3pt}
【消息】\par
〈按句编号的消息文本：[1] …… [2] ……〉\par
\end{quote}}
